\section{A conditional nonlinear tail transfer}
\label{app:transfer}

For completeness, we state the precise measure-theoretic wrapper around the
deterministic estimate \eqref{eq:det-transfer}.  It records which missing
input is needed to use the exact Gaussian law under a nonlinear gauge measure.

\begin{proposition}[Conditional transfer]
Let $\mu$ be a target measure, $\nu$ a Gaussian reference measure, and
$\mathcal G$ an event on which a common coupling supplies a centered surrogate
$x$ and logarithmic staple vectors $v$ with
\[
 \norm{C_mx-v}_{\ell^2}\le\sqrt m\,\varepsilon.
\]
Assume also that for every measurable $A$,
\[
 \mu(A\cap\mathcal G)\le e^L\nu(A\cap\mathcal G),
 \qquad \mu(\mathcal G^c)\le q_{\mathrm{exc}}.
\]
Then
\[
 \mu(D\ge m\eta^2)
 \le e^L\nu\bigl(Q\ge m(\eta-\varepsilon)_+^2\bigr)
 +q_{\mathrm{exc}}.
\]
\end{proposition}

\begin{proof}
On $\mathcal G$, the reverse triangle inequality gives
$\sqrt Q\ge\sqrt D-\sqrt m\varepsilon$.  Hence the target event restricted to
$\mathcal G$ is contained in the displayed Gaussian event.  Apply the density
comparison on $\mathcal G$ and add the exceptional probability.
\end{proof}

This proposition is deliberately conditional.  The finite-block spectrum
settles the reference tail exactly, while the constants $L$, $\varepsilon$,
and $q_{\mathrm{exc}}$ require nonlinear, state-uniform analysis not supplied
by the certificate.
